\chapter{Fecal Leukocytes (White Blood Cells in Stool)}

\section*{What Is This Test?}
A fecal leukocyte test detects white blood cells (neutrophils) in stool, indicating intestinal inflammation. CAUSES of inflammatory diarrhea include invasive bacterial infection (Shigella, Salmonella, Campylobacter, EHEC, C.\ difficile) and inflammatory bowel disease. Non-inflammatory causes (viral, parasitic, toxigenic bacterial, osmotic, secretory) lack substantial fecal leukocytes. The test has largely been superseded by fecal calprotectin and lactoferrin, which are more sensitive and quantitative, but stool leukocyte microscopy remains a rapid bedside or office test in resource-limited settings.

\section*{Alternative Names}
Fecal Leukocytes; Stool White Blood Cells; Stool Smear for WBC; Fecal Wright's or Methylene Blue Stain

\section*{Principle}
A thin smear of stool stained with methylene blue or Wright's stain is examined microscopically at 400x for polymorphonuclear leukocytes. Reporting semi-quantitative: 0 to many per HPF. Modern kits use lactoferrin latex agglutination (Leuko-Test) as a surrogate marker of PMNs but are mostly replaced by fecal calprotectin immunoassay.

\section*{History}
Fecal leukocyte microscopy was first described in the 1940s as a way to distinguish bacillary dysentery from amoebic and non-inflammatory diarrhea. Methylene blue staining was standardized in the 1960s.

\section*{Why This Test Is Ordered}
\begin{itemize}
  \item Differentiate inflammatory from non-inflammatory acute diarrhea
  \item Evaluate bloody or mucoid diarrhea
  \item Suspected invasive bacterial enteritis (Shigella, Salmonella, Campylobacter)
  \item Suspected C.\ difficile colitis and pseudomembranous colitis
  \item Workup of chronic diarrhea to distinguish IBD from irritable bowel syndrome
\end{itemize}

\section*{Primary vs Secondary Test}
Primary in resource-limited settings; otherwise secondary to fecal calprotectin for IBD screening and stool culture for infectious diarrhea.

\section*{How This Test Is Performed}
A small stool sample is smeared on a slide, fixed, stained with methylene blue or Wright's stain, and examined at 400x. A fresh specimen is preferred; samples $>$ 2--4 hours old may have degraded leukocytes.

\section*{What Preparation Is Needed}
None; barium and oily laxatives interfere with microscopy. Antidiarrheal agents do not affect findings during active disease.

\section*{Contraindications and Precautions}
Low sensitivity (15--30\% false negatives) limits utility: many patients with invasive bacterial diarrhea have few leukocytes in stool; fecal leukocytes falsely low in treated or resolving colitis. Specimen degradation over 2 hours; sample preservation crucial.

\section*{Risks}
None; stool collection only.

\section*{What Happens After the Test}
Positive result prompts stool culture and C.\ difficile testing. Negative result does not exclude invasive disease; clinical judgment supersedes.

\section*{Understanding Results}

\subsection*{Below Normal}
\begin{itemize}
  \item Non-inflammatory diarrhea: viral (rotavirus, norovirus), toxigenic bacterial (cholera, ETEC), parasitic (Giardia)
  \item IBS or functional disorders
  \item Invasive bacterial diarrhea with low leukocyte shedding: false negative
\end{itemize}

\subsection*{Above Normal}
\begin{itemize}
  \item Invasive bacterial enteritis: Shigella, Salmonella, Campylobacter, EIEC, EHEC
  \item Clostridioides difficile colitis (often many leukocytes)
  \item Ulcerative colitis or Crohn's disease (active)
  \item Ischemic colitis, radiation colitis
  \item Rare: Yersinia, Aeromonas
\end{itemize}

\section*{Normal Results}
\begin{itemize}
  \item \textbf{Fecal leukocytes per HPF:} 0 to rare (occasional)
  \item \textbf{Positive result:} $>$ 3--5 PMN/HPF or ``moderate'' to ``many'' on smear
  \item \textbf{Clinical context:} True positive is sensitive when stools are grossly bloody, mucoid, or during acute flare
\end{itemize}

\section*{Follow-up Tests}
\begin{itemize}
  \item \textbf{Stool culture (Salmonella, Shigella, Campylobacter, E.\ coli O157)}: Identifies invasive pathogens
  \item \textbf{C.\ difficile PCR and toxin assay}: If antibiotic-associated diarrhea suspected
  \item \textbf{ Stool O\&P examination}: If parasitic etiology considered (Giardia, Entamoeba)
  \item \textbf{Fecal calprotectin or lactoferrin:} More sensitive markers of intestinal inflammation for IBD workup
  \item \textbf{CBC with differential and CRP}: Systemic inflammation
  \item \textbf{Stool PCR multiplex panel}: Comprehensive molecular diagnosis of enteric pathogens
  \item \textbf{Colonoscopy with biopsy}: If IBD suspected based on chronic symptoms
\end{itemize}

\section*{References}
\begin{enumerate}
  \item Thielman NM, Guerrant RL. Acute infectious diarrhea. N Engl J Med. 2004;350(1):38-47.
  \item Guerrant RL, et al. Practice guidelines for the management of infectious diarrhea. Clin Infect Dis. 2001;32(3):357-364.
  \item Mosli M, et al. Fecal calprotectin as a diagnostic biomarker for IBD. Am J Gastroenterol. 2015;110(8):1099-1106.
  \item Harris JC, et al. Fecal leukocytes in infectious diarrhea. Ann Intern Med. 1972;76(4):561-565.
\end{enumerate}

\clearpage
\section*{Regulatory Status and Authorization Date}
Regulatory status applies to the specific assay, instrument, manufacturer, and intended use rather than to the test name alone. No single approval date can be assigned to this test category. For a commercial assay, verify the current product in the relevant regulator's database: the U.S. Food and Drug Administration (FDA) clearance, approval, or authorization; India's Central Drugs Standard Control Organisation (CDSCO) licence or permission; the European Union In Vitro Diagnostic Regulation (EU IVDR) CE certificate issued through the applicable conformity-assessment route; or the United Kingdom Medicines and Healthcare products Regulatory Agency (MHRA) registration and UKCA/CE status. Laboratory-developed tests may not have an individual FDA, CDSCO, EU IVDR, or MHRA approval date.
\clearpage
